\section{Background and related work}\label{sec:background}

\subsection{Kinetic model discrimination}
For a library $\mathcal{M}$ of candidate mechanisms $m_1,\dots,m_K$ with parameter vectors $\boldsymbol{\theta}_m$, the conventional workflow calibrates every candidate by weighted least squares and compares the calibrated candidates through an information criterion. With Gaussian measurement errors of known standard deviation the Bayesian information criterion is \citep{Schwarz1978}
\begin{equation}
\BIC_m = \chi^2_m(\hat{\boldsymbol{\theta}}_m) + k_m \ln n ,
\label{eq:bic_bg}
\end{equation}
where $\chi^2_m$ is the minimised sum of squared standardised residuals, $k_m$ the number of estimated parameters and $n$ the number of measurements. Differences $\Delta_m = \BIC_m - \min_{m'}\BIC_{m'}$ convert into weights $w_m \propto \exp(-\Delta_m/2)$ that can be read as approximate posterior model probabilities under equal priors \citep{Burnham2002}. When the error variances are unknown they can be estimated jointly with the parameters; for responses with separate variances this leads to a log-determinant or sum-of-log criterion \citep{BoxDraper1965,Bard1974}. Information criteria rank candidates relative to one another; whether the best candidate is adequate in absolute terms is a separate question, answered by comparing its residuals with the measurement error. Model-based design of experiments chooses operating conditions that maximise the expected divergence between rival models \citep{HunterReiner1965,BoxHill1967,BuzziFerraris1983} or the precision of their parameters \citep{Franceschini2008}. Parameter identifiability, assessed for instance through profile likelihoods \citep{Raue2009}, determines how much of a calibrated model can be interpreted.

\subsection{Hybrid models and neural ordinary differential equations}
Hybrid models replace the poorly known part of a first-principles model with a data-driven component, most commonly a neural network that predicts specific rates inside the mass balances \citep{Psichogios1992,Thompson1994,Oliveira2004}. The known structure improves data efficiency and extrapolation relative to black-box models, which has made hybrid models popular in bioprocess engineering \citep{vonStosch2014,Narayanan2020,Sansana2021,Bradley2022}. When the hybrid model is an ODE, its trajectories can be differentiated with respect to the network weights through the numerical integrator, so the network can be trained on the measured states directly rather than on rates estimated by numerical differentiation; this is the neural ODE \citep{Chen2018,Kidger2021}, also termed a universal differential equation when the network is embedded in an otherwise mechanistic model \citep{Rackauckas2020}. Physics-informed neural networks follow a different route, representing the trajectory itself by a network and penalising the residual of the governing equations \citep{Raissi2019}; they have been used to estimate the kinetic parameters of a given mechanism in fermentations \citep{Bangi2022}. Neither route, on its own, tells the modeller which mechanism the culture follows.

\subsection{Equation discovery and symbolic regression}
Sparse identification of nonlinear dynamics (SINDy) writes the time derivative of the state as a sparse linear combination of library functions and selects the active terms by sparsity-promoting regression \citep{Brunton2016}. Rate laws of biochemical kinetics are rational functions, which are not sparse in a polynomial library; implicit formulations recover them by multiplying through by the denominator, so that $\mu\,(K_S + S) = \mu_{\max} S$ becomes linear in its coefficients \citep{Mangan2016,Kaheman2020}. Reactive SINDy and chemical reaction neural networks adapt these ideas to reaction networks \citep{Hoffmann2019,JiDeng2021}, and genetic-programming symbolic regression searches over unrestricted expression trees \citep{Schmidt2009,Cranmer2023}. Most of these methods need time derivatives of the state, which are difficult to obtain from a dozen noisy offline samples. Learning the unknown term with a neural ODE first and applying sparse regression to its output avoids numerical differentiation \citep{Rackauckas2020}. The workflow presented here builds on that idea, but treats the outcome of the regression as a set of hypotheses to be tested against the measurements alongside the established mechanisms, constrains the hypotheses to physically admissible rate laws, and closes the loop by adding accepted laws to the library.
